\begin{proposition}
	\label{prop:upperbound-automorphisms}
	If \( A \subseteq \Aut(L) \) is a subgroup, then \( [L:L'] \leq \# A \).
\end{proposition}

\begin{proof}
	We argue by contradiction. Suppose that \( [L:L'] > \# A \). Say we have \( \# A = n \), and that
	\( u_1,\ldots, u_{n+1} \in L \) being linearly independent over \( L' \). Consider the system of
	linear equations:
	\[
		\sum_{i=1}^{n+1} \underbrace{\sigma(u_i)}_{\in L}  x_i = 0 \quad \forall \; \sigma \in A
	\]
	since \# equations < \# variables, then there exists \( a_1,\ldots, a_{n+1} \in L \) which are not
	all \( 0 \), s.t.
	\[
		\sum_{i=1}^{n+1} a_i \sigma (u_i) = 0 \quad \forall \; \sigma  \in A
	\]
	Let \( q \leq n+1 \) be minimal, such that we have
	\[
		\sum_{i=1}^q b_i \sigma (v_i) = 0 \quad \forall \; \sigma \in A \text{ for some } v_1,\ldots,
		v_q \in L \text{ being linearly independent over } L'
	\]
	such that not all \( b_i \) are \( 0 \). By minimality of \( q \), all \( b_i \ne 0 \), one can
	replace all \( b_i \) by \( b_i / b_1 \) and thus may assume \( b_1 = 1 \). We want to then get a
	shorter such relation:
	\begin{itemize}
		\item If \( q = 1 \): we have \( \sigma(v_1) = 0 \), leading to contradiction with the fact that
		      \( \sigma \) is injective and \( v_1 \ne 0 \).
		\item May assume that \( q \geq 2 \), let \( \sigma = \Id \) in the this case, we thus have:
		      \[
			      \sum_{i=1}^q b_i v_i = 0
		      \]
		      As the \( v_i \) are linearly independent over \( L' \implies \exists \; i \geq 2 \), s.t.
		      \( b_i \not\in L' \), otherwise it will be \( 0 \). After reordering, may assume such \( i = 2 \). We
		      have that \( b_2 \not\in L' \implies \exists \; \tau \in A \), s.t. \( \tau (b_2) \ne b_2 \), now apply
		      \( \tau \) to the equation, we have:
		      \[
			      \sum_{i=1}^q \tau(b_i) \tau \sigma (v_i) = 0 \quad \forall \; \sigma \in A
		      \]
		      Since \( A \) is a subgroup, see that \( \tau \sigma \in A \), by letting \( \sigma = \tau
		      \sigma\), we have:
		      \[
			      \sum_{i=1}^q b_i \tau \sigma (v_i) = 0 \quad \forall \; \sigma \in A
		      \]
		      We then subtract the two things, and get
		      \[
			      \begin{aligned}
				      \sum_{i=1}^q (\tau (b_i) - b_i) \tau \sigma (v_i) & = 0 \quad \forall \; \sigma  \in A
				      \\
				      \implies \sum_{i=1}^q \underbrace{(\tau (b_i) - b_i)}_{\substack{=0 \text{ if } i = 1  \\ \ne 0 \text{
				      if } i = 2}} \sigma (v_i)                         & = 0 \forall \; \sigma \in A
			      \end{aligned}
		      \]
		      thus we get a shorter relation, contradicting to the minimality of \( q \) \( \lightning \).
	\end{itemize}
\end{proof}

\begin{theorem}[\textbf{Fundemental Theorem of Galois Theory}]
	Given in the finite extension setting, let \( K \hookrightarrow L \) be a finite Galois extension:
	\begin{enumerate}
		\item These two maps describe above gives \textcolor{red}{mutual inverse order-reversing bijection} between
		      subextension of \( K \hookrightarrow L \) and subgroups of \( G \).
		\item The subgroup \( H \) of \( G \) is a normal subgroup of \( G \) iff \( K \hookrightarrow
		      L^H \) is a normal extension. In this case, we have an isormorphism:
		      \[
			      \qo{G}{H} \cong G\bace{\qo{L^H}{K}}
		      \]
	\end{enumerate}
\end{theorem}

\begin{proof}[\textbf{Proof of Fundemental Theorem of Galois Theory}]
	\leavevmode
	\begin{enumerate}
		\item If \( K \hookrightarrow L' \hookrightarrow L \) being a subextension, with \( [L:L'] = m \), since
		      \( \qo{L}{L'} \) is Galois extension, see that \( \abs{G(\qo{L}{L'})} = m \). By \textbf{Proposition} \ref{prop:upperbound-automorphisms}
		      and \ref{prop:lowerbound-automorphisms}, see that
		      \[
			      [L:L^{G(\qo{L}{L'})}] = \abs{G(\qo{L}{L'})} = m
		      \]
		      and we have the trivial inclusion:
		      \[
			      \underbrace{L' \subseteq \underbrace{L^{G(\qo{L}{L'})} \subseteq L}}_{\text{same degree}} \implies L^{G(\qo{L'}{L})} = L'
		      \]
		      and similarly: if \( H \leq G(\qo{L}{K}) \), let \( \abs{H} = q \), by
		      \textbf{Proposition} \ref{prop:upperbound-automorphisms} and
		      \ref{prop:lowerbound-automorphisms}, see that
		      \[
			      [L:L^H] = m \implies \abs{G\bace{\qo{L}{L^H}}} = m
		      \]
		      and trivially we have:
		      \[
			      H\leq G\bace{\qo{L}{L^H}}
		      \]
		      and they have the same order, thus \( H = G(\qo{L}{L^H}) \).
		\item Now suppose \( K \hookrightarrow L' \hookrightarrow L \) be a subextension, s.t. \(
		      \qo{L'}{K} \) is normal extension. If \( \sigma \in G(\qo{L}{K}) \), since \( \qo{L'}{K} \) is normal, then
		      \[
			      \sigma(L') \subseteq L' \implies \sigma |_{L'} \in G\bace{\qo{L'}{K}}
		      \]
		      we thus get a natrual map as follow, which is in fact a group homomorphism:
		      \[
			      \begin{aligned}
				      vp: G\bace{\qo{L}{K}} & \to G\bace{\qo{L'}{K}} \\
				      \sigma                & \rsa \sigma |_{L'}
			      \end{aligned}
		      \]
		      One can see that \( \ker(\vp) = G(\qo{L}{L'}) \implies G(\qo{L}{L'}) \) is a normal subgroup
		      of \( G \), thus we have:
		      \[
			      \qo{G}{\ker(\vp)} \xrightarrow{\sim} \im(\vp) \xhookrightarrow{\sim} G\bace{\qo{L'}{K}}
		      \]
		      by that the order of \( \qo{G}{\ker(\vp)}) \) is given by:
		      \[
			      \abs{\qo{G}{\ker(\vp)}} = \frac{\abs{G}}{\abs{G\bace{\qo{L}{L'}}}} = \frac{[L:K]}{[L:L']}
			      = [L':K] = \abs{G\bace{\qo{L'}{K}}}
		      \]
		      thus the isomorphism holds:
		      \[
			      \qo{G}{G\bace{\qo{L}{L'}}} \cong G\bace{\qo{L'}{K}}
		      \]
		      Suppose now \( K \hookrightarrow L' \hookrightarrow L \) and let \( H = G(\qo{L}{L'}) \)
		      and that \( H \trianglelefteq G \). We want to see that \( \qo{L'}{K} \) is a normal extension.
		      \[
			      \begin{tikzcd}
				      K \arrow[r,hook] \arrow[rd, hook] & L' \arrow[r,hook] \arrow[rrd, "\alpha"] & L
				      \arrow[dashed, rd, "\tilde{\alpha}"] & \\
				      & L' \arrow[r,hook] & L \arrow[r,hook] &
				      \overline{L}
			      \end{tikzcd}
		      \]
		      We need to show: \( \forall \) morphism of \( K \)-extension \( \alpha: L' \to \overline{L} \), we have
		      that \( \alpha(L') \subseteq L' \). Since \( \qo{L}{L'} \) is algebraic, we can get an
		      extension of \( \alpha \) by \( \tilde{\alpha} \). Since \( \qo{L}{K} \) is normal, then
		      \( \tilde{\alpha}(L) \subseteq L \implies \sigma\coloneqq  \tilde{\alpha}|_{L} \in G(\qo{L}{K})  \). We
		      need to show: \( \sigma(L') \subseteq L' \). \textbf{By 1}, we already know: \( L' = L^H \), hence we need to
		      show: if \( u \in L^H \implies \sigma(u) \in L^H \). Which is equivalent to show:
		      \[
			      \begin{aligned}
				      \forall \; \tau \in H,              \quad  \tau \sigma (u) & = \sigma (u) \\
				      \iff (\sigma ^{-1}\tau \sigma )(u)                         & = u
			      \end{aligned}
		      \]
		      with \( \sigma^{-1} \tau \sigma \in H \) since \( H \trianglelefteq G \). Thus \( u\in L^H
		      \implies \sigma^{-1} \tau \sigma(u) = u\).
	\end{enumerate}
\end{proof}

\begin{eg}
	Let \( \qo{L}{\QQ} \) be the splitting field of \( x^3 - 2 \), where
	\[
		L = \QQ(\sqrt[3]{2}, \omega) \quad \omega^2 + \omega + 1 = 0
	\]
	\begin{remark}
		In general, if \( \qo{L}{K} \) is a splitting field of a polynomial \( f \) (\( \deg(f) = n \)),
		let \( \text{char}(K) = 0 \), and define:
		\[
			K \hookrightarrow L = K(\alpha_1,\ldots, \alpha_n)
		\]
		where \( \alpha_1,\ldots, \alpha_n \) are roots of \( f \). Given \( \sigma \in G(\qo{L}{K})
		\implies \forall \; i, \; \sigma(\alpha_i) = \alpha_j \) for some \( j \). Thus we have the map:
		\[
			\begin{aligned}
				G(\qo{L}{K}) & \to \fS_n                                                 \\
				\sigma       & \rsa (i \to j \text{ if } \sigma (\alpha _i) = \alpha _j)
			\end{aligned}
		\]
		this is a group homomorphism, and injective since \( L = K(\alpha_1,\ldots, \alpha_n) \). In
		general this is very far from surjective.
	\end{remark}
	In this case, we in fact have a surjective and thus bijective homomorphism:
	\[
		G\bace{\qo{L}{\QQ}} \hookrightarrow \fS_3
	\]
	this has to be an isomorphism, since otherwise, \( \qo{L}{\QQ} \) would have no non-trivial
	subextension, and thus \( \fS_3 \) have no non-trivial subgroup, but it has.

	Alternatively, we have:
	\[
		\underbrace{\QQ \hookrightarrow}_{\deg = 3} \QQ(\sqrt[3]{2}) \underbrace{\hookrightarrow
			L}_{\deg = 2}
	\]
	with the latter one given by the fact that \( \omega \not \in \QQ(\sqrt[3]{2}) \) and \( \omega^2
	+ \omega + 1 = 0\). Thus \( \qo{L}{\QQ} \) has degree to be \( 6 \), the corresponding group
	should either be \( \qo{\ZZ}{6\ZZ} \) or \( \qo{\ZZ}{2\ZZ} \times \qo{\ZZ}{3\ZZ} \) or \( \fS_3
	\). But then since \( \qo{\QQ(\sqrt[3]{2})}{\QQ} \) is \textbf{not normal}, thus \( G(\qo{L}{\QQ}) \ne \)
	abelian group, thus
	\[
		G\bace{\qo{L}{\QQ}} \cong \fS_3
	\]
	One can view in this way:
	\( \fS_3 \) has an index \( 2 \) subgroup, \( 3 \) subgroups of order \( 2 \), which \textbf{\textcolor{blue}{correponds to}}
	\[
		\begin{tikzcd}
			\QQ \arrow[r,hook] \arrow[rd, "\deg = 3"',hook] \arrow[rdd, hook,"\deg=3"'] \arrow[rddd, hook,
				"\deg = 3"'] & \QQ(\omega) \arrow[r,hook] & L \\
			& \QQ(\sqrt[3]{2}) \arrow[ru,hook, "\deg=2"'] & \\
			& \QQ(\sqrt[3]{2}\omega) \arrow[ruu, hook, "\deg=2"'] & \\
			& \QQ(\sqrt[3]{2}\omega^2) \arrow[ruuu, hook, "\deg=2"'] &
		\end{tikzcd}
	\]
\end{eg}

